\section{Conclusion}
\label{sec:conclusion}

We show that when a modern instruction-tuned LLM is asked to lie, the resulting text is strongly and detectably different from truthful text at the surface level.
A simple logistic regression classifier using 21 linguistic features achieves 90.0\% accuracy at distinguishing truthful from deceptive responses, driven primarily by response brevity, reduced hedging, and increased certainty language.
This contradicts the earlier finding of \citet{schuster2019limitations} that language models show no stylistic differences when generating deceptive text, and suggests that RLHF training has created a detectable ``truthful mode'' that models depart from when instructed to lie.

Our results have implications for AI safety: deceptive LLM outputs may be detectable from surface features alone, without requiring access to model internals.
The reversal of the human hedging pattern---where LLMs become \emph{more} assertive when lying, not less---highlights a fundamental difference between human and machine deception that deserves further study.

Future work should validate these findings across multiple model families, test robustness to adversarial prompt engineering designed to eliminate stylistic tells, and investigate whether the lying style persists in longer-form generation.
The relationship between internal truth representations \citep{burger2024truth} and the surface-level signals we identify also warrants exploration: do models whose internals encode truth more strongly show larger surface-level shifts when forced to lie?
